% GENERATED FILE. Edit canonical lessons, then run npm run generate:books.
% canonical-source-hash: fnv1a64:21bff5f0c757b8c5
% canonical-lessons: AR-C527-fallah, AR-C527-bustani, AR-C527-rabiha, AR-C527-aflasa, AR-C527-iddakhara, AR-R527-review-520-527

\chapter{Farmers, profits and savings}
\label{ch:ar-farmers-profits-and-savings}
\hlchaptermodality{527}

\begin{chapteropening}
\textbf{By the end of this chapter:} \emph{I can name a peasant farmer and a gardener, and say that someone made a profit, went bankrupt or put money aside.}

Complete the last lesson of chapter 527: I can name a peasant farmer and a gardener, and say that someone made a profit, went bankrupt or put money aside.
\end{chapteropening}

\section[\texorpdfstring{\ar{فلاح} (fallāḥ)}{فلاح (fallāḥ)}]{\ar{فلاح} (fallāḥ) — a peasant farmer}

\label{lesson:AR-C527-fallah}



\begin{quote}
Before the new one: say the Arabic for a builder, then the Arabic for a cobbler.
\end{quote}

\subsection*{You'll want to know: \ar{فلاح}}
\textbf{\ar{فلاح}} — \emph{fallāḥ} — \textquotedblleft{}a peasant farmer\textquotedblright{}.

\begin{quote}
\ar{يزرع} \ar{الفلاح} \ar{القمح} \ar{في} \ar{الربيع}.
\end{quote}

\begin{quote}
\emph{yazraʿu al-fallāḥ al-qamḥ fī ar-rabīʿ.}
\end{quote}

\begin{quote}
The farmer sows wheat in the spring.
\end{quote}

\begin{cousinweb}
\textbf{\ar{فلاح}} (\emph{fallāḥ}) is built on the native Arabic root \emph{f-l-ḥ}, \textquotedblleft{}to split, to till the soil; to prosper\textquotedblright{} (Lane). A fallāḥ is one who splits the earth to till it.
\end{cousinweb}

\begin{grammarlens}[title={What you've built}]
One more word for describing what you lived and felt.
\end{grammarlens}

\subsection*{Your turn}
\begin{itemize}
\raggedright
  \item \emph{Say it:} \emph{fallāḥ}
  \item \emph{Say it:} \emph{fallāḥ}, once more
  \item \emph{Say it:} \emph{yazraʿu al-fallāḥ al-qamḥ fī ar-rabīʿ.} — The farmer sows wheat in the spring.
  \item \emph{Recall:} say the Arabic for a builder, then the Arabic for a cobbler, then say \emph{fallāḥ} again
\end{itemize}

\begin{tcolorbox}[breakable,colback=teal!4,colframe=teal!35!black,title={Before you move on}]
What does \ar{فلاح} mean? (\textquotedblleft{}A peasant farmer\textquotedblright{}.) Say it once more.
\end{tcolorbox}
\section[\texorpdfstring{\ar{بستاني} (bustānī)}{بستاني (bustānī)}]{\ar{بستاني} (bustānī) — a gardener}

\label{lesson:AR-C527-bustani}



\begin{quote}
Before the new one: say the Arabic for a cobbler, then the Arabic for a peasant farmer.
\end{quote}

\subsection*{You'll want to know: \ar{بستاني}}
\textbf{\ar{بستاني}} — \emph{bustānī} — \textquotedblleft{}a gardener\textquotedblright{}.

\begin{quote}
\ar{يسقي} \ar{البستاني} \ar{الورد} \ar{كل} \ar{صباح}.
\end{quote}

\begin{quote}
\emph{yasqī al-bustānī al-ward kull ṣabāḥ.}
\end{quote}

\begin{quote}
The gardener waters the roses every morning.
\end{quote}

\begin{cousinweb}
\textbf{\ar{بستاني}} (\emph{bustānī}) is built on \emph{bustān}, \textquotedblleft{}an orchard\textquotedblright{}, which you know, with the ending \emph{-ī}: the orchard man. \emph{Bustān} is usually traced to Persian.
\end{cousinweb}

\begin{grammarlens}[title={What you've built}]
One more word for describing what you lived and felt.
\end{grammarlens}

\subsection*{Your turn}
\begin{itemize}
\raggedright
  \item \emph{Say it:} \emph{bustānī}
  \item \emph{Say it:} \emph{bustānī}, once more
  \item \emph{Say it:} \emph{yasqī al-bustānī al-ward kull ṣabāḥ.} — The gardener waters the roses every morning.
  \item \emph{Recall:} say the Arabic for a cobbler, then the Arabic for a peasant farmer, then say \emph{bustānī} again
\end{itemize}

\begin{tcolorbox}[breakable,colback=teal!4,colframe=teal!35!black,title={Before you move on}]
What does \ar{بستاني} mean? (\textquotedblleft{}A gardener\textquotedblright{}.) Say it once more.
\end{tcolorbox}
\section[\texorpdfstring{\ar{ربح} (rabiḥa)}{ربح (rabiḥa)}]{\ar{ربح} (rabiḥa) — to make a profit}

\label{lesson:AR-C527-rabiha}



\begin{quote}
Before the new one: say the Arabic for a peasant farmer, then the Arabic for a gardener.
\end{quote}

\subsection*{You'll want to know: \ar{ربح}}
\textbf{\ar{ربح}} — \emph{rabiḥa} — \textquotedblleft{}to make a profit\textquotedblright{}.

\begin{quote}
\ar{ربح} \ar{التاجر} \ar{كثيرا} \ar{هذه} \ar{السنة}.
\end{quote}

\begin{quote}
\emph{rabiḥa at-tājir kathīran hādhihi as-sana.}
\end{quote}

\begin{quote}
The merchant made a big profit this year.
\end{quote}

Present: \ar{يربح}, \emph{yarbaḥ}.

In a game it also means to win.

\begin{cousinweb}
\textbf{\ar{ربح}} (\emph{rabiḥa}) is built on the native Arabic root \emph{r-b-ḥ}, \textquotedblleft{}to gain in trade\textquotedblright{} (Lane). The noun ribḥ is a profit.
\end{cousinweb}

\begin{grammarlens}[title={What you've built}]
One more word for describing what you lived and felt.
\end{grammarlens}

\subsection*{Your turn}
\begin{itemize}
\raggedright
  \item \emph{Say it:} \emph{rabiḥa}
  \item \emph{Say it:} \emph{rabiḥa}, once more
  \item \emph{Say it:} \emph{rabiḥa at-tājir kathīran hādhihi as-sana.} — The merchant made a big profit this year.
  \item \emph{Recall:} say the Arabic for a peasant farmer, then the Arabic for a gardener, then say \emph{rabiḥa} again
\end{itemize}

\begin{tcolorbox}[breakable,colback=teal!4,colframe=teal!35!black,title={Before you move on}]
What does \ar{ربح} mean? (\textquotedblleft{}To make a profit\textquotedblright{}.) Say it once more.
\end{tcolorbox}
\section[\texorpdfstring{\ar{أفلس} (aflasa)}{أفلس (aflasa)}]{\ar{أفلس} (aflasa) — to go bankrupt}

\label{lesson:AR-C527-aflasa}



\begin{quote}
Before the new one: say the Arabic for a gardener, then the Arabic for to make a profit.
\end{quote}

\subsection*{You'll want to know: \ar{أفلس}}
\textbf{\ar{أفلس}} — \emph{aflasa} — \textquotedblleft{}to go bankrupt\textquotedblright{}.

\begin{quote}
\ar{أفلست} \ar{الشركة} \ar{بعد} \ar{سنتين}.
\end{quote}

\begin{quote}
\emph{aflasat ash-sharika baʿda sanatayn.}
\end{quote}

\begin{quote}
The company went bankrupt after two years.
\end{quote}

Present: \ar{يفلس}, \emph{yuflis}.

\begin{cousinweb}
\textbf{\ar{أفلس}} (\emph{aflasa}) is built on the native Arabic root \emph{f-l-s} and is usually traced to \emph{fals} (plural \emph{fulūs}), a small copper coin: to be left with coppers and no silver.
\end{cousinweb}

\begin{grammarlens}[title={What you've built}]
One more word for describing what you lived and felt.
\end{grammarlens}

\subsection*{Your turn}
\begin{itemize}
\raggedright
  \item \emph{Say it:} \emph{aflasa}
  \item \emph{Say it:} \emph{aflasa}, once more
  \item \emph{Say it:} \emph{aflasat ash-sharika baʿda sanatayn.} — The company went bankrupt after two years.
  \item \emph{Recall:} say the Arabic for a gardener, then the Arabic for to make a profit, then say \emph{aflasa} again
\end{itemize}

\begin{tcolorbox}[breakable,colback=teal!4,colframe=teal!35!black,title={Before you move on}]
What does \ar{أفلس} mean? (\textquotedblleft{}To go bankrupt\textquotedblright{}.) Say it once more.
\end{tcolorbox}
\section[\texorpdfstring{\ar{ادخر} (iddakhara)}{ادخر (iddakhara)}]{\ar{ادخر} (iddakhara) — to put aside (savings)}

\label{lesson:AR-C527-iddakhara}



\begin{quote}
Before the new one: say the Arabic for to make a profit, then the Arabic for to go bankrupt.
\end{quote}

\subsection*{You'll want to know: \ar{ادخر}}
\textbf{\ar{ادخر}} — \emph{iddakhara} — \textquotedblleft{}to put aside (savings)\textquotedblright{}.

\begin{quote}
\ar{ادخرت} \ar{أختي} \ar{مالا} \ar{لشراء} \ar{سيارة}.
\end{quote}

\begin{quote}
\emph{iddakharat ukhtī mālan li-shirāʾ sayyāra.}
\end{quote}

\begin{quote}
My sister put money aside to buy a car.
\end{quote}

Present: \ar{يدخر}, \emph{yaddakhir}.

\ar{وفر}, to save, which you know, is the everyday verb.

\begin{cousinweb}
\textbf{\ar{ادخر}} (\emph{iddakhara}) is built on the native Arabic root \emph{dh-kh-r}, \textquotedblleft{}to store up\textquotedblright{} (Lane). Form VIII iddakhara, with dh turning into d, is to store a thing for later.
\end{cousinweb}

\begin{grammarlens}[title={What you've built}]
One more word for describing what you lived and felt.
\end{grammarlens}

\subsection*{Your turn}
\begin{itemize}
\raggedright
  \item \emph{Say it:} \emph{iddakhara}
  \item \emph{Say it:} \emph{iddakhara}, once more
  \item \emph{Say it:} \emph{iddakharat ukhtī mālan li-shirāʾ sayyāra.} — My sister put money aside to buy a car.
  \item \emph{Recall:} say the Arabic for to make a profit, then the Arabic for to go bankrupt, then say \emph{iddakhara} again
\end{itemize}

\begin{tcolorbox}[breakable,colback=teal!4,colframe=teal!35!black,title={Before you move on}]
What does \ar{ادخر} mean? (\textquotedblleft{}To put aside (savings)\textquotedblright{}.) Say it once more.
\end{tcolorbox}
\section[(dialogue)]{Words from chapters 520–527 — a review}

\label{lesson:AR-R527-review-520-527}



\begin{quote}
Say the Arabic for to get used to and a shift.
\end{quote}

\subsection*{Your turn}
Say these aloud:

\begin{itemize}
\raggedright
  \item to get used to, a shift, time off, a task, to carry out
  \item a wage, a bonus, a promotion, to promote, to get tired
  \item to set to work, to toil, a resignation, retired, unemployed
  \item to make plans, to complete, to devote oneself, an evaluation, production
  \item to hold talks, to work hard, to be on duty, to be promoted, headquarters
  \item to produce, to design, a designer, to program, a programmer
  \item a mechanic, a building contractor, a house painter, a builder, a cobbler
  \item a peasant farmer, a gardener, to make a profit, to go bankrupt, to put aside
\end{itemize}

\begin{tcolorbox}[breakable,colback=teal!4,colframe=teal!35!black,title={Before you move on}]
To get used to? (\textbf{\ar{اعتاد}}, \emph{iʿtāda}.) A shift? (\textbf{\ar{مناوبة}}, \emph{munāwaba}.) Time off? (\textbf{\ar{إجازة}}, \emph{ijāza}.) A task? (\textbf{\ar{مهمة}}, \emph{mahamma}.) To carry out? (\textbf{\ar{نفذ}}, \emph{naffadha}.) A wage? (\textbf{\ar{أجر}}, \emph{ajr}.) A bonus? (\textbf{\ar{مكافأة}}, \emph{mukāfaʾa}.) A promotion? (\textbf{\ar{ترقية}}, \emph{tarqiya}.) To promote? (\textbf{\ar{رقى}}, \emph{raqqā}.) To get tired? (\textbf{\ar{تعب}}, \emph{taʿiba}.) To set to work? (\textbf{\ar{باشر}}, \emph{bāshara}.) To toil? (\textbf{\ar{كدح}}, \emph{kadaḥa}.) A resignation? (\textbf{\ar{استقالة}}, \emph{istiqāla}.) Retired? (\textbf{\ar{متقاعد}}, \emph{mutaqāʿid}.) Unemployed? (\textbf{\ar{عاطل}}, \emph{ʿāṭil}.) To make plans? (\textbf{\ar{خطط}}, \emph{khaṭṭaṭa}.) To complete? (\textbf{\ar{أتم}}, \emph{atamma}.) To devote oneself? (\textbf{\ar{تفرغ}}, \emph{tafarragha}.) An evaluation? (\textbf{\ar{تقييم}}, \emph{taqyīm}.) Production? (\textbf{\ar{إنتاج}}, \emph{intāj}.) To hold talks? (\textbf{\ar{فاوض}}, \emph{fāwaḍa}.) To work hard? (\textbf{\ar{اجتهد}}, \emph{ijtahada}.) To be on duty? (\textbf{\ar{ناوب}}, \emph{nāwaba}.) To be promoted? (\textbf{\ar{ترقى}}, \emph{taraqqā}.) Headquarters? (\textbf{\ar{مقر}}, \emph{maqarr}.) To produce? (\textbf{\ar{أنتج}}, \emph{antaja}.) To design? (\textbf{\ar{صمم}}, \emph{ṣammama}.) A designer? (\textbf{\ar{مصمم}}, \emph{muṣammim}.) To program? (\textbf{\ar{برمج}}, \emph{barmaja}.) A programmer? (\textbf{\ar{مبرمج}}, \emph{mubarmij}.) A mechanic? (\textbf{\ar{ميكانيكي}}, \emph{mīkānīkī}.) A building contractor? (\textbf{\ar{مقاول}}, \emph{muqāwil}.) A house painter? (\textbf{\ar{دهان}}, \emph{dahhān}.) A builder? (\textbf{\ar{بناء}}, \emph{bannāʾ}.) A cobbler? (\textbf{\ar{إسكافي}}, \emph{iskāfī}.) A peasant farmer? (\textbf{\ar{فلاح}}, \emph{fallāḥ}.) A gardener? (\textbf{\ar{بستاني}}, \emph{bustānī}.) To make a profit? (\textbf{\ar{ربح}}, \emph{rabiḥa}.) To go bankrupt? (\textbf{\ar{أفلس}}, \emph{aflasa}.) To put aside? (\textbf{\ar{ادخر}}, \emph{iddakhara}.)
\end{tcolorbox}
